\paragraph{Annotated}
   with this type you can put comments for example in the following example x muist be positive
\begin{minted}{python}
      from typing import Annotated
      x: Annotated[int, "positive"] = 42
   \end{minted}
   with fast API, this example say that the integer myst be between 0 and 100
\begin{minted}{python}
      from typing import Annotated
      from fastapi import Query
      
      def get_items(limit: Annotated[int, Query(gt=0, lt=100)]):
          return {"limit": limit}
   \end{minted}

\paragraph{Optional}
   
   When the type is a class object or None

\begin{minted}{python}
      ->Optional[CommentedMap]
   \end{minted}
   
   The above code says the return is either CommentedMap or None

\paragraph{Generic}
\begin{minted}{python}
      from typing import Generic, TypeVar, Protocol
      T = TypeVar("T")
      class Unit(Protocol, Generic[T]):
          def __iter__(self) -> Iterator[T]: ...
          def __len__(self) -> int: ...
   \end{minted}



\paragraph{Mapping}
\begin{minted}{python}
       data: Mapping[str, Any] = field(default_factory=dict)
   \end{minted}


\paragraph{ModuleType}
   For modules (ie a file with many functions or classes)  that you want to refer to all of the at once such az \pyl{import matplotlib.pyplot as plt} the only way is to use
   \begin{minted}{python}
   from types import ModuleType

   plt: ModuleType
   \end{minted}

\paragraph{Final}